\section{Partial volume: the BrainWeb phantom}\label{sec:pv}

The forward model assumes one tissue per voxel. Real voxels at 1--1.2\,mm often contain
several, and the cortex, two to three voxels thick, borders CSF almost everywhere. This
section measures the consequences with the BrainWeb anatomical model and the unchanged
estimators of \cref{sec:ml,sec:magnitude,sec:paper}, and explains them from the structure of
the model.

\subsection{The mixed-voxel signal}

If a voxel contains compartments $c$ with volume fractions $f_c$, proton densities $\mathrm{PD}_c$ and
relaxation parameters $\theta_c$, and shares the location-dependent quantities $\Bone$, $\dw$, $\phi_0$
and the receive sensitivity $C$, its signal is
\begin{equation}
  S^{\rm mix}_{ijk}=C\sum_cf_c\,\mathrm{PD}_c\;s_{ijk}(\theta_c),\qquad
  s_{ijk}(\theta_c)=v_{i,c}\,G_k(u_{i,c},E_{2,c},\dw)\,\E^{-t/T_{2,c}}\E^{-R'_{2,c}|t+k\TR|}\E^{i\Phi},
  \label{eq:mix}
\end{equation}
by \cref{thm:uv} applied to each compartment. A single-compartment estimator returns the
$\hat\theta$ whose model $S(\hat\theta)$ is closest to $S^{\rm mix}$ (in the sense of its own cost function).
Two questions follow: how far $S^{\rm mix}$ lies from the set of single-compartment signals (the
model manifold), which determines bias and detectability, and where on the manifold the
closest point lies, which determines the reported values.

\paragraph{Nearby compartments: a second-order effect.} If the compartments are close in
parameter space, write $\theta_c=\bar\theta+\delta_c$ with $\bar\theta$ the weighted mean (weights
$w_c=f_c\mathrm{PD}_c/\sum f\mathrm{PD}$, in log coordinates). Then
\begin{equation}
  \frac{S^{\rm mix}}{C\sum f\mathrm{PD}}=\sum_cw_cS(\theta_c)
  =S(\bar\theta)+\tfrac12\sum_cw_c\,\delta_c^{\mathsf T}\nabla^2S(\bar\theta)\,\delta_c+O(|\delta|^3),
  \label{eq:mix2}
\end{equation}
because the first-order terms cancel ($\sum w_c\delta_c=0$). The mixture is a single-compartment
signal at the weighted mean, up to a second-order term whose size is set by the curvature of
the model and the spread of the compartments. White and grey matter are close
($\Tone$ 850 vs 1350\,ms, $\Ttwo$ 66 vs 90\,ms): a WM/GM mixture is fitted to within $0.01\sigma$ per
measurement at SNR 1000 and reported at the weighted mean (median deviation $-0.013$ in
$\log\Tone$, \cref{tab:bwclean}). Such voxels are harmless: they look like an intermediate tissue,
which is what they are.

\paragraph{CSF: a first-order mismatch.} CSF is far from tissue ($\Tone=4000$\,ms, $\Ttwo=1500$\,ms),
so \eqref{eq:mix2} does not apply, and its signal contribution is strongly pathway- and
scan-dependent. Per unit volume, CSF produces $1.4$--$1.7$ times the grey-matter signal in the
FID of scan 1, but $2.6$--$3.4$ times in the echo pathways of scan 1 and $6.5$--$7.8$ times in
pathway $+1$ of scan 2 (proton-density ratio $1.22$): long $\Ttwo$ preserves the higher-order
configuration states that the echo pathways sample, and the large second pulse feeds them. A
voxel with $10\%$ CSF by volume therefore takes about $13\%$ of its FID but up to about $46\%$ of
its scan-2, $k=+1$ signal from CSF. No single $(\Tone,\Ttwo,\Ttwos)$ reproduces this pattern: the mixture
lies off the manifold, the residual is first order in the CSF fraction, and the fitted values
depend on which measurements an estimator weights most.

\paragraph{Consequences.} (i) Apparent parameters are not any simple average: relative to the
PD-weighted mean they are biased in a direction and by an amount that depend on the
parameter and the estimator (\cref{fig:bwcurves,tab:bwclean}). (ii) Different estimators no
longer agree, because they project the same off-manifold point differently: complex ML
weights all 36 measurements, the analytic inverse estimates $\Ttwo$ from scan 1 alone and $\Bone$
from closed-form ratios. (iii) The shared $\Bone$ absorbs part of the mismatch: the two scans see
the CSF admixture differently, and an estimator that ties them by a common $\Bone$ may shift
$\Bone$ to reconcile them. (iv) The residual grows only linearly with the CSF fraction while the
noise is fixed, so detectability depends on SNR: the noise-free cost $\|r\|^2$ makes the
residual $\chi^2$ non-central with parameter $\lambda=\|r\|^2/\sigma^2$, which grows as SNR$^2$.
(v) Uncertainty maps describe noise, not model error, so their coverage must drop wherever
the bias is comparable to the noise.

\subsection{Phantom and protocol}

We use the BrainWeb normal-brain anatomical model~\cite{collins1998,kwan1999,cocosco1997}: fuzzy
tissue-membership volumes at 1\,mm isotropic resolution ($181\times217\times181$). An axial slice through
the lateral ventricles (index 95, Talairach $z=23$\,mm) is zero-padded to $256\times256$. Voxels whose
discrete label is CSF, grey or white matter are kept (\,18\,668 voxels\,); fat, skin and skull are
excluded, since fat would require a chemical-shift model. Each compartment takes the tissue
values of \cref{tab:tissues} with a $3\%$ smooth texture, and the fields $\Bone$, $\dw$, $C$, $\phi_0$ are
those of \cref{sec:phantom}, rescaled to the head. Signals follow \eqref{eq:mix}. Of the
brain voxels, 13\,258 are pure (one tissue $\ge95\%$), 2818 are WM/GM mixtures and 2592 contain
5--95\% CSF (\cref{fig:bwphantom}).

Three cases are reconstructed with the estimators unchanged (four starts plus the spatial
restart for the two least-squares fits): \emph{noise-free with partial volume} (model error
alone), \emph{SNR 1000 with partial volume}, and a \emph{control} at SNR 1000 in which each
voxel is assigned entirely to its dominant tissue (same anatomy, no partial volume). For a
mixed voxel we report errors against the PD-weighted geometric mean
$\theta_{\rm eff}=\exp(\sum_cw_c\log\theta_c)$, and show the dominant-tissue value alongside, since a mixed
voxel has no single true value.

\begin{figure}[!htbp]
\centering
\includegraphics[width=0.92\linewidth]{brainweb_phantom.png}
\caption{BrainWeb slice: tissue fractions, the PD-weighted effective $\Tone$ and $\Ttwo$, and voxel
categories (pure; WM/GM mixed; containing $\ge5\%$ CSF).}
\label{fig:bwphantom}
\end{figure}

\subsection{Results}

\begin{table}[!htbp]
\centering
\caption{BrainWeb, noise-free: median of $\log(\hat\theta/\theta_{\rm eff})$ per voxel category, where $\theta_{\rm eff}$ is the PD-weighted geometric mean of the compartment values ($\Bone$: relative to the true field). Last columns: RMS mismatch residual per measurement in units of $\sigma$ at SNR 1000, and the probability that a $\chi^2$ test (99.9\,\%) flags the voxel at SNR 1000 and 3000.}
\label{tab:bwclean}
\footnotesize
\setlength{\tabcolsep}{3pt}
\begin{tabular}{lrcccccccc}
\toprule
 & & \multicolumn{3}{c}{Complex ML} & \multicolumn{3}{c}{Analytic inverse} & & \\
\cmidrule(lr){3-5}\cmidrule(lr){6-8}
Category & voxels & $\Tone$ & $\Ttwo$ & $\Bone$ & $\Tone$ & $\Ttwo$ & $\Bone$ & RMS/$\sigma$ & flagged 1000 / 3000\\
\midrule
pure CSF & 1823 & +0.000 & +0.000 & +0.000 & +0.000 & +0.000 & +0.000 & 0.00 & 0.1\,\% / 0.1\,\%\\
pure GM & 3931 & +0.000 & +0.000 & +0.000 & +0.000 & +0.000 & +0.000 & 0.00 & 0.1\,\% / 0.1\,\%\\
pure WM & 7504 & +0.000 & +0.000 & +0.000 & +0.000 & +0.000 & +0.000 & 0.00 & 0.1\,\% / 0.1\,\%\\
WM/GM mixed & 2818 & -0.013 & -0.009 & +0.000 & -0.017 & -0.006 & +0.000 & 0.01 & 0.1\,\% / 0.1\,\%\\
CSF 5-25\% & 1069 & -0.133 & -0.038 & +0.005 & -0.042 & -0.142 & -0.007 & 0.35 & 3.4\,\% / 59.2\,\%\\
CSF 25-50\% & 580 & -0.324 & -0.103 & +0.009 & -0.113 & -0.428 & -0.032 & 0.61 & 12.4\,\% / 95.7\,\%\\
CSF 50-95\% & 943 & -0.251 & -0.113 & +0.002 & -0.029 & -0.503 & -0.012 & 0.29 & 2.2\,\% / 48.1\,\%\\
\bottomrule
\end{tabular}
\end{table}

\begin{table}[!htbp]
\centering
\caption{BrainWeb at SNR$(\Mzero)=1000$, complex ML: SD of $\log(\hat\Tone/T_{1,\rm eff})$ with partial volume and in the control without it, coverage of $95\%$ uncertainty intervals, and the fraction of voxels flagged by the $\chi^2$ test.}
\label{tab:bwnoisy}
\footnotesize
\begin{tabular}{lccccc}
\toprule
Category & $\Tone$ SD (PV) & $\Tone$ SD (control) & median (PV) & coverage (PV) & flagged (PV)\\
\midrule
pure CSF & 0.077 & 0.077 & -0.003 & 94.7\,\% & 0.0\,\%\\
pure GM & 0.060 & 0.060 & +0.000 & 94.6\,\% & 0.2\,\%\\
pure WM & 0.084 & 0.084 & +0.000 & 94.9\,\% & 0.1\,\%\\
WM/GM mixed & 0.057 & 0.049 & -0.012 & 93.8\,\% & 0.1\,\%\\
CSF 5-25\% & 0.085 & 0.063 & -0.138 & 21.9\,\% & 4.0\,\%\\
CSF 25-50\% & 0.058 & 0.052 & -0.318 & 0.5\,\% & 15.2\,\%\\
CSF 50-95\% & 0.139 & 0.079 & -0.243 & 26.8\,\% & 3.1\,\%\\
\bottomrule
\end{tabular}
\end{table}


\begin{figure}[!htbp]
\centering
\includegraphics[width=\linewidth]{brainweb_mixture_curves.png}
\caption{Noise-free GM/CSF voxels ($<5\%$ WM): values returned by the single-compartment
estimators versus the CSF volume fraction, with the PD-weighted mean and the grey-matter value
(nominal, without texture). Right: $\Bone$ error.}
\label{fig:bwcurves}
\end{figure}

\begin{figure}[!htbp]
\centering
\includegraphics[width=\linewidth]{brainweb_maps.png}
\caption{Complex ML at SNR$(\Mzero)=1000$. From left: estimate with partial volume; log-ratio to the
dominant tissue's value; log-ratio to the PD-weighted mean; the same anatomy without partial
volume (control). Rows: $\Tone$ and $\Ttwo$.}
\label{fig:bwmaps}
\end{figure}

\paragraph{Pure tissue and WM/GM mixtures.} Pure voxels are reconstructed without bias by all
estimators, and complex ML has the same spread and $94.6$--$94.9\%$ coverage as in the control
(\cref{tab:bwnoisy}). WM/GM mixtures are fitted to $0.01\sigma$ per measurement, reported at the
PD-weighted mean to within $-0.013$ in $\log\Tone$ and $-0.009$ in $\log\Ttwo$ for complex ML, and their
coverage ($93.8\%$) is close to nominal: the second-order picture of \eqref{eq:mix2} holds.

\paragraph{CSF-mixed voxels: biased, confidently.} In the 2592 voxels with $5$--$95\%$ CSF all
estimators are biased by far more than the noise (\cref{tab:bwclean,fig:bwcurves,fig:bwmaps}).
For complex ML the deviation of $\log\Tone$ from the PD-weighted mean is $-0.13$, $-0.32$ and $-0.24$
in the three CSF bins; $\Ttwo$ deviates by $-0.04$ to $-0.11$ from the mean but lies far above the
grey-matter value (\cref{fig:bwcurves}: at $10\%$ CSF the apparent $\Ttwo$ is about $40\%$ above that of
grey matter). Against the dominant tissue, the cortex therefore appears with elevated $\Ttwo$ and,
where CSF dominates, depressed $\Tone$ (\cref{fig:bwmaps}). The noise spread in these voxels is
similar to the control (\cref{tab:bwnoisy}), so the uncertainty maps report small intervals
around a biased value: coverage of $95\%$ intervals drops to $21.9\%$, $0.5\%$ and $26.8\%$, against
$94$--$95.5\%$ in the control. The residual does not reveal the problem at this SNR: the mismatch
is $0.29$--$0.61\sigma$ per measurement, and a $\chi^2$ test at $99.9\%$ flags $3$--$15\%$ of these voxels at
SNR 1000 (predicted $2.2$--$12.4\%$ from the noise-free residual), rising to $48$--$96\%$ only at SNR
3000.

\paragraph{The estimators disagree.} Complex ML and magnitude least squares return the same
values (both fit all measurements), but the analytic inverse does not: its $\Ttwo$ deviates by
$-0.14$ to $-0.50$ from the mean (it uses scan 1 only), and in $19.5$--$38\%$ of the CSF-mixed voxels
the mismatch moves its flip-angle equation onto a different root, with $\Bone$ about $40\%$ too low
and $\Tone$ of $4$--$9$\,s even without noise (\cref{fig:bwcurves}, right). Complex ML's $\Bone$ stays within
$+4\%$ in every mixed voxel; its error is confined to the relaxation times.

\subsection{What partial volume does to the model and the estimators}

\begin{enumerate}
\item \textbf{Partial volume is a model error, not noise.} It biases every single-compartment
  estimator, does not average out with more data, and is invisible to the CRLB, which assumes
  the model is correct. Efficiency comparisons (\cref{sec:ml}) remain valid in pure voxels,
  where they were made.
\item \textbf{Its size is governed by distance on the model manifold.} Mixtures of nearby
  tissues are, to second order, single-compartment signals at the weighted mean \eqref{eq:mix2};
  they are harmless and should be read as intermediate tissue. Mixtures with CSF are far from the
  manifold because CSF's long $\Ttwo$ inflates its contribution to the echo pathways and to the
  second scan, up to eight times its volume share: the reported values then depend on the
  measurements an estimator weights most, and no estimator returns a well-defined average.
\item \textbf{It is hard to detect per voxel.} The mismatch residual grows linearly with the
  CSF fraction while the noise is fixed, so a $\chi^2$ test detects partial volume only at high SNR
  or after averaging. Per-voxel uncertainty maps become falsely confident where it acts; they
  should be reported together with a partial-volume indicator, e.g.\ the distance to a tissue
  boundary from a segmentation, or the fit of a mixture model.
\item \textbf{The shared $\Bone$ matters.} Because $\Bone$ couples the two scans, an estimator that uses the
  scans in separate closed-form steps (the analytic inverse) can resolve the mismatch by a large
  $\Bone$ error, whereas the joint fit keeps $\Bone$ nearly unbiased and pushes the error into the
  relaxation times. For the two-stage design this favours estimating $\Bone$ from robustly smoothed
  maps, which the polynomial fit already provides: partial-volume errors are local and are
  rejected as outliers.
\item \textbf{Remedy.} A two-compartment model,
  $S=C\,[f\,\mathrm{PD}_as(\theta_a)+(1-f)\,\mathrm{PD}_bs(\theta_b)]$ with shared $\Bone$, $\dw$, $\phi_0$, is the natural
  correction. Fitting both compartments freely is ill-conditioned (13 parameters, near-collinear
  signals), but with class parameters fixed from pure voxels or a population estimate (the
  hierarchical model of \cref{sec:future}) the signal is linear in the fractions given $\Bone$, $\dw$,
  $\phi_0$, and variable projection reduces the fit to a small nonlinear problem. CSF fractions should
  then be well determined, because CSF's signature is distinct, whereas WM/GM fractions will be
  poorly determined---and, by item 2, need not be. The implicit Jacobian, the CRLB machinery and
  the FID-sign test carry over unchanged (every compartment's FID has the sign of $\sin\alpha$).
\end{enumerate}
